\section{LLM Agent 之前的历史}

\subsection{基于规则的 Text Agent}

\begin{knowledgebox}{ELIZA（1966）}
最早的聊天机器人之一，仅用少量规则（不断追问或复述用户所说）就能给人以"智能"的印象。然而，基于规则的方法具有严重的局限性：
\begin{itemize}
    \item 高度领域特定（domain-specific）
    \item 每个新领域需要重新编写规则
    \item 无法扩展到复杂开放域
\end{itemize}
\end{knowledgebox}

\subsection{基于强化学习的 Text Agent}

在 LLM 出现之前，另一个流行范式是用 RL 构建 Text Agent（如文字游戏 agent）：

\begin{itemize}
    \item 将文本观测和动作视为视频游戏中的像素和键盘输入
    \item 使用 RL 优化奖励信号
    \item \textbf{局限}：领域特定、需要精确的奖励信号、训练开销大
\end{itemize}

\subsection{LLM 的革命性变化}

LLM 仅通过在海量文本上训练 next-token prediction，就能在推理时通过 prompting 解决各种新任务。这种\textbf{通用性}和\textbf{少样本学习能力}使其成为构建 Agent 的理想``大脑''。

\subsection{本章小结}

从基于规则的 symbolic AI agent 到基于 RL 的 deep agent，再到基于 LLM 的 reasoning agent，Agent 的内部表示从符号逻辑到神经嵌入再到自然语言，一步步走向更灵活、更通用。

